\documentclass[12pt]{article}
\usepackage[utf8]{inputenc}
\usepackage[T1]{fontenc}
\usepackage[margin=1in]{geometry}
\usepackage{setspace}
\usepackage{xcolor}
\usepackage[colorlinks=true, urlcolor=blue!60!black, linkcolor=blue!60!black]{hyperref}
\setlength{\parindent}{2em}
\setlength{\parskip}{1.25em}
\renewcommand{\baselinestretch}{1.0}
\title{\textit{Puerto Rico v. Branstad}, 483 U.S. 219 (1987)}
\author{Case Brief}
\date{}
\begin{document}
\maketitle

\section{Case Metadata}
\begin{itemize}
  \item \textbf{Case Name:} \emph{Puerto Rico v. Branstad}, \href{https://supreme.justia.com/cases/federal/us/483/219/}{483 U.S. 219 (1987)}
  \item \textbf{Docket No.:} 85-2116 \quad \textbf{Argued:} March 30, 1987 \quad \textbf{Decided:} June 23, 1987
  \item \textbf{Court:} Supreme Court of the United States (on writ of certiorari to the U.S. Court of Appeals for the Eighth Circuit, 787 F.2d 423 (8th Cir. 1986))
  \item \textbf{Briefing Book Placement:} Chapter 4: Federalism and State Sovereignty
  \item \textbf{Vote:} 9--0 judgment reversing the Eighth Circuit (6--3 as to Part II-B's constitutional holding)
  \item \textbf{Opinion of the Court:} Justice Marshall (joined in full by Chief Justice Rehnquist and Justices Brennan, White, Blackmun, and Stevens; joined in Parts I, II-A, II-C, and III by Justices Powell and O'Connor; joined in part by Justice Scalia)
  \item \textbf{Concurring Opinions:} Justice O'Connor (joined by Justice Powell), concurring in part and concurring in the judgment; Justice Scalia, concurring in part and concurring in the judgment
  \item \textbf{Dissenting Opinions:} None
  \item \textbf{Key Constitutional and Statutory Provisions:} Extradition Clause, U.S. Const. art. IV, \S~2, cl. 2; Tenth Amendment; Federal Extradition Act, 18 U.S.C. \S~3182 (originally Act of Feb. 12, 1793, 1 Stat. 302)
  \item \textbf{Primary Source Links:} \href{https://supreme.justia.com/cases/federal/us/483/219/}{Justia U.S. Supreme Court Center}; \href{https://www.law.cornell.edu/constitution-conan/article-4/section-2/clause-2/overview-of-extradition-interstate-rendition-clause}{Cornell LII Constitution Annotated (Extradition Clause Overview)}
\end{itemize}

\section{Facts and Procedural History}
On January 25, 1981, Ronald Calder, an Iowa native working as a civilian Federal Aviation Administration air traffic controller in San Juan, Puerto Rico, struck two pedestrians with his automobile in a grocery store parking lot in Aguadilla, Puerto Rico, following an altercation with Antonio de Jesus Gonzalez. De Jesus Gonzalez was injured, and his wife, Army Villalba---then eight months pregnant---was killed along with her unborn child. Two sworn statements given to police, including one from De Jesus Gonzalez and one from an eyewitness, reported that after striking the couple, Calder backed his car over Villalba's prostrate body two or three times (\href{https://supreme.justia.com/cases/federal/us/483/219/}{483 U.S. at 221}).

Arrested and initially released on \$5,000 bail, Calder was arraigned on February 4, 1981, in a District Court of the Commonwealth of Puerto Rico on charges of first-degree murder and attempted murder. He failed to appear at a preliminary hearing on March 4, 1981 (prompting a bail increase to \$50,000) and again failed to appear on April 13, 1981, at which point he was declared a fugitive from justice and bail was raised to \$300,000. Believing Calder had fled to his family's home in Iowa, Puerto Rico police alerted Iowa authorities; Calder surrendered in Polk County, Iowa, on April 24, 1981, posted a \$20,000 bond, and was released (\href{https://supreme.justia.com/cases/federal/us/483/219/}{483 U.S. at 221--22}).

On May 15, 1981, the Governor of Puerto Rico transmitted a formal extradition requisition to the Governor of Iowa, accompanied by the arrest warrant, fugitive resolution, charging documents, and three sworn witness statements. At an extradition hearing conducted by the Iowa Governor's counsel on June 17, 1981, Calder's attorney testified that ``a white American man . . . could not receive a fair trial in the Commonwealth of Puerto Rico,'' and Calder asserted that witnesses in Puerto Rican courts had been ``bought'' (\href{https://supreme.justia.com/cases/federal/us/483/219/}{483 U.S. at 222}). After negotiations to reduce the charges failed, Iowa Governor Robert Ray formally denied extradition on December 28, 1981, citing the absence of a ``change to a more realistic charge.'' A renewed extradition request to Ray's successor, Governor Terry Branstad, was likewise refused (\href{https://supreme.justia.com/cases/federal/us/483/219/}{483 U.S. at 222}).

After the Supreme Court denied Puerto Rico leave to file an original-jurisdiction bill of complaint (\emph{Puerto Rico v. Iowa}, 464 U.S. 1034 (1984)), Puerto Rico filed suit on February 15, 1984, in the U.S. District Court for the Southern District of Iowa against Governor Branstad and Iowa, seeking declaratory relief and a writ of mandamus compelling Branstad to perform the ministerial duty of extradition under the Extradition Clause and 18 U.S.C. \S~3182. Although Iowa stipulated that Puerto Rico's extradition papers fully satisfied the statute, the District Court dismissed the suit as barred by \href{https://supreme.justia.com/cases/federal/us/65/66/}{\emph{Kentucky v. Dennison}, 65 U.S. (24 How.) 66 (1861)}. The Eighth Circuit ``[r]eluctantly'' affirmed (787 F.2d 423, 424 (8th Cir. 1986)), and the Supreme Court granted certiorari (479 U.S. 811 (1986)).

\section{Issue(s) Presented}
\begin{enumerate}
  \item Does \href{https://supreme.justia.com/cases/federal/us/65/66/}{\emph{Kentucky v. Dennison}} remain good law insofar as it held that federal courts lack constitutional authority to issue a writ of mandamus compelling a state governor to perform the mandatory, ministerial duty to extradite a fugitive under the Extradition Clause and the Extradition Act?
  \item May the Commonwealth of Puerto Rico---though not a ``State'' within the literal text of Article IV, \S~2, cl. 2---invoke federal judicial power to compel a state governor to extradite a fugitive under the Federal Extradition Act, 18 U.S.C. \S~3182?
\end{enumerate}

\section{Holding}
\begin{enumerate}
  \item \textbf{No.} Overruling \href{https://supreme.justia.com/cases/federal/us/65/66/}{\emph{Kentucky v. Dennison}}, the Supreme Court held that federal courts have authority to issue a writ of mandamus compelling the governor of an asylum State to perform the mandatory, ministerial duty to deliver up a fugitive upon proper demand under the Extradition Clause and the Extradition Act (\href{https://supreme.justia.com/cases/federal/us/483/219/}{483 U.S. at 227--30}).
  \item \textbf{Yes.} Without deciding whether the Extradition Clause applies directly to Puerto Rico, the Court held that the Extradition Act (18 U.S.C. \S~3182) applies to Puerto Rico as a Territory/Commonwealth, empowering it to seek federal mandamus relief against a state governor (\href{https://supreme.justia.com/cases/federal/us/483/219/}{483 U.S. at 229--30}).
\end{enumerate}

\section{Reasoning}
Writing for the Court, Justice Marshall dismantled \emph{Dennison}'s no-enforcement rule in three steps.

\textbf{1. Historical Context of \emph{Dennison} (Part II-A).} Marshall situated \href{https://supreme.justia.com/cases/federal/us/65/66/}{\emph{Kentucky v. Dennison}} in its immediate pre-Civil War crisis. Decided on March 14, 1861---after Deep South states had seceded and weeks before Fort Sumter---Chief Justice Taney's opinion arose when ``the practical power of the Federal Government [was] at its lowest ebb since the adoption of the Constitution'' (\href{https://supreme.justia.com/cases/federal/us/483/219/}{483 U.S. at 225}). Taney had actually held that the Extradition Clause imposed an absolute, ``merely ministerial'' duty on Ohio's Governor to surrender a free Black man indicted in Kentucky for helping an enslaved woman escape, yet Taney then neutered that duty as merely ``moral,'' declaring that ``the Federal Government, under the Constitution, has no power to impose on a State officer, as such, any duty whatever, and compel him to perform it'' (65 U.S. at 107).

\textbf{2. Obsolescence of Coequal-Sovereigns Federalism and Tenth Amendment Analysis (Part II-B).} Marshall reaffirmed \emph{Dennison}'s first proposition---that the Extradition Clause is mandatory and leaves asylum-state officers no discretion, consistent with \href{https://supreme.justia.com/cases/federal/us/439/282/}{\emph{Michigan v. Doran}, 439 U.S. 282 (1978)}, and \href{https://supreme.justia.com/cases/federal/us/482/400/}{\emph{California v. Superior Court}, 482 U.S. 400 (1987)}---while repudiating its second. More than a century of post-Civil War constitutional jurisprudence had rejected Taney's premise that States and the Federal Government are in all respects coequal sovereigns. Federal courts routinely enforce federal constitutional duties against state officers through mandamus (\emph{Board of Liquidation v. McComb}, 92 U.S. 531 (1876)) and prospective injunctive relief (\href{https://supreme.justia.com/cases/federal/us/209/123/}{\emph{Ex parte Young}, 209 U.S. 123 (1908)}; \emph{Brown v. Board of Education}; \emph{Cooper v. Aaron}). Because the extradition obligation is ministerial, explicit, and imposed directly on the States by the Constitution itself, compelling compliance avoids ongoing federal supervision of discretionary state governance and requires no Tenth Amendment balancing (\href{https://supreme.justia.com/cases/federal/us/483/219/}{483 U.S. at 228}). Marshall also rejected Iowa's defense of an ``executive common law'' of gubernatorial discretion: entrenched noncompliance cannot rewrite a clear constitutional command (\href{https://supreme.justia.com/cases/federal/us/483/219/}{483 U.S. at 228--29}).

\textbf{3. Puerto Rico's Statutory Right Under 18 U.S.C. \S~3182 (Part II-C).} Although Article IV, \S~2, cl. 2 refers only to ``States,'' the Court avoided deciding whether the Clause directly covers Puerto Rico because 18 U.S.C. \S~3182 expressly encompasses any demanding ``Territory.'' Under \emph{New York ex rel. Kopel v. Bingham}, 211 U.S. 468 (1909), Puerto Rico could demand extradition under the Act as a Territory, and its 1950s transition to Commonwealth status did not strip it of that authority. Because federal courts may enforce federal statutory duties against state officials (\emph{Maine v. Thiboutot}, 448 U.S. 1 (1980)), Puerto Rico could maintain its mandamus action under \S~3182 (\href{https://supreme.justia.com/cases/federal/us/483/219/}{483 U.S. at 229--30}).

\section{Concurring and Dissenting Opinions}
\begin{itemize}
  \item \textbf{Justice O'Connor (joined by Justice Powell), concurring in part and concurring in the judgment:} Joined Parts I, II-A, II-C, and III, but declined to join Part II-B. Invoking constitutional avoidance (\emph{Ashwander v. TVA}, 297 U.S. 288, 347 (1936)), she reasoned that because 18 U.S.C. \S~3182 fully resolved Puerto Rico's right to mandamus, the Court's broader pronouncements regarding the Extradition Clause and the Tenth Amendment were unnecessary (\href{https://supreme.justia.com/cases/federal/us/483/219/}{483 U.S. at 230--31}).
  \item \textbf{Justice Scalia, concurring in part and concurring in the judgment:} Concurred in the result and solely in the portions of the opinion applying 18 U.S.C. \S~3182 and overruling \emph{Dennison} as an interpretation of the 1793 statute's predecessor. Presaging his later anti-commandeering jurisprudence, he pointedly noted ``that no party before us has asserted the lack of power of Congress to require extradition from a State to a Territory'' (\href{https://supreme.justia.com/cases/federal/us/483/219/}{483 U.S. at 231}).
  \item \textbf{Dissents:} None.
\end{itemize}

\section{Key Quotations}
\begin{quote}
  ``We reaffirm the conclusion that the commands of the Extradition Clause are mandatory, and afford no discretion to the executive officers or courts of the asylum State.''
  --- Justice Marshall, \href{https://supreme.justia.com/cases/federal/us/483/219/}{483 U.S. at 227}.
\end{quote}
\begin{quote}
  ``Because the duty is directly imposed upon the States by the Constitution itself, there can be no need to weigh the performance of the federal obligation against the powers reserved to the States under the Tenth Amendment.''
  --- Justice Marshall, \href{https://supreme.justia.com/cases/federal/us/483/219/}{483 U.S. at 228}.
\end{quote}
\begin{quote}
  ``Long continuation of decisional law or administrative practice incompatible with the requirements of the Constitution cannot overcome our responsibility to enforce those requirements.''
  --- Justice Marshall, \href{https://supreme.justia.com/cases/federal/us/483/219/}{483 U.S. at 229}.
\end{quote}
\begin{quote}
  ``\emph{Kentucky v. Dennison} is the product of another time. The conception of the relation between the States and the Federal Government there announced is fundamentally incompatible with more than a century of constitutional development. Yet this decision has stood while the world of which it was a part has passed away. We conclude that it may stand no longer.''
  --- Justice Marshall, \href{https://supreme.justia.com/cases/federal/us/483/219/}{483 U.S. at 230} (also quoted in \href{https://www.law.cornell.edu/constitution-conan/article-4/section-2/clause-2/overview-of-extradition-interstate-rendition-clause}{Cornell LII Constitution Annotated}).
\end{quote}

\section{Doctrinal Significance and Modern Legacy (Through 2026)}
\emph{Branstad} occupies a pivotal hinge point in horizontal and vertical federalism. First, it eliminated the 126-year anomaly of an unenforceable constitutional text, cementing interstate rendition as a summary, federally enforceable obligation. Eleven years later, in \emph{New Mexico ex rel. Ortiz v. Reed}, 524 U.S. 151, 155 (1998) (cited in \href{https://www.law.cornell.edu/constitution-conan/article-4/section-2/clause-2/overview-of-extradition-interstate-rendition-clause}{Cornell LII Constitution Annotated}), the Court unanimously applied \emph{Branstad} to forbid an asylum state from refusing extradition based on a fugitive's claim of duress or fear of unfair treatment in the demanding state.

Second, \emph{Branstad} marks the outer boundary of the modern \textbf{anti-commandeering doctrine}. Just five to ten years after \emph{Branstad}, the Court revived Tenth Amendment limits on federal compulsion of state legislatures (\href{https://supreme.justia.com/cases/federal/us/505/144/}{\emph{New York v. United States}, 505 U.S. 144 (1992)}), state executive officers (\href{https://supreme.justia.com/cases/federal/us/521/898/}{\emph{Printz v. United States}, 521 U.S. 898 (1997)}), and state regulatory prohibitions (\href{https://supreme.justia.com/cases/federal/us/584/16-476/}{\emph{Murphy v. NCAA}, 584 U.S. 453 (2018)}). Justice Marshall's careful Tenth Amendment carve-out in \emph{Branstad}---that extradition is a duty ``directly imposed upon the States by the Constitution itself''---reconciles \emph{Branstad} with \emph{Printz}: Congress may not commandeer state executives to enforce ordinary Article I regulatory programs, but federal courts may compel state governors to execute an express Article IV constitutional mandate.

Third, through 2025--2026, \emph{Branstad} has returned to the center of sharp interstate constitutional conflicts. Following \emph{Dobbs v. Jackson Women's Health Organization} (2022), states banning abortion have sought to prosecute out-of-state telehealth providers, while ``shield law'' states have barred gubernatorial extradition. In February 2025, Louisiana Governor Jeff Landry demanded the extradition of New York physician Dr. Margaret Carpenter for prescribing medication abortion pills into Louisiana, and New York Governor Kathy Hochul refused under New York's shield law (\href{https://www.brookings.edu/articles/louisiana-raises-the-stakes-in-the-state-abortion-wars/}{Brookings Institution, May 8, 2025}). Crucially, \emph{Branstad} does \emph{not} compel extradition in such telehealth cases because Article IV, \S~2, cl. 2 and 18 U.S.C. \S~3182 apply only to a ``fugitive'' who was physically present in the demanding state and ``fled from Justice''; rendering non-fugitives rests on discretionary state statutes (such as Uniform Criminal Extradition Act \S~6). By contrast, in August 2026, when Texas Governor Greg Abbott refused Minnesota's extradition demand for a suspect who indisputably fled Minnesota, commentators noted that \emph{Branstad} and \emph{Ortiz} squarely foreclose a governor's policy-based refusal to surrender an actual fugitive (\href{https://www.stevevladeck.com/p/bonus-244-minnesota-v-abbott}{Steve Vladeck, Aug. 20, 2026}).

\section{Pedagogical Critiques and Exam Hypotheticals}
\textbf{Pedagogical Critiques.} First, notice the analytical mismatch flagged by Justices O'Connor and Scalia: Part II-B invokes a duty ``directly imposed upon the States by the Constitution itself'' to bypass the Tenth Amendment, yet Part II-C declines to hold that the Extradition Clause applies to Puerto Rico and rests the mandamus remedy solely on a federal statute (18 U.S.C. \S~3182). If Congress's power to extend extradition to Territories is purely statutory rather than textually mandated by Article IV, \emph{Printz}'s anti-commandeering rule creates genuine tension with \emph{Branstad}'s Territory holding---precisely the question Scalia reserved in his one-paragraph concurrence. Second, \emph{Branstad} highlights a recurring structural tension in horizontal federalism: mandatory, no-discretion rendition promotes interstate harmony and prevents sanctuaries, yet strips asylum-state governors of any humanitarian safety valve when demanding states pursue politically or morally contested prosecutions (the exact dilemma that animated antebellum free-state governors in \emph{Dennison}).

\textbf{Exam Hypothetical and ``Opposite Side'' Contrast.}
\begin{itemize}
  \item \textbf{Scenario A (Mandatory Extradition Under \emph{Branstad}):} State X indicts Dr. A for performing a procedure illegal in State X while Dr. A was physically operating a clinic in State X. Dr. A then moves to State Y, whose Governor refuses State X's facially valid extradition demand because State Y's legislature passed a statute prohibiting extradition for that offense. \emph{Analysis:} State X will prevail in federal court via mandamus under \emph{Branstad} and \emph{Doran}. Dr. A was physically present in State X and fled; State Y's shield law is preempted by Article IV, \S~2, cl. 2 and 18 U.S.C. \S~3182, and the Tenth Amendment offers no defense.
  \item \textbf{Scenario B (``Opposite Side'' Contrast---Non-Fugitive Telehealth \& Anti-Commandeering):} Now suppose Dr. B has never set foot in State X; from an office in State Y, Dr. B prescribes medication via telehealth to a patient in State X. State X indicts Dr. B and demands extradition; State Y's Governor refuses under State Y's shield law. Congress then enacts a new federal statute under the Commerce Clause requiring every state governor to arrest and extradite any non-fugitive accused of a cross-border telehealth crime. \emph{Analysis:} On the existing Extradition Clause and \S~3182, State Y wins because Dr. B never ``fled from Justice'' in State X, so \emph{Branstad}'s mandatory duty is never triggered. Moreover, the hypothetical Commerce Clause statute would likely be invalidated under \href{https://supreme.justia.com/cases/federal/us/521/898/}{\emph{Printz}} and \href{https://supreme.justia.com/cases/federal/us/584/16-476/}{\emph{Murphy}}: unlike \emph{Branstad}, where the Constitution itself directly imposed the duty on state governors, Congress cannot use Article I to commandeer state executives into enforcing a non-constitutional extradition mandate.
\end{itemize}

\end{document}
